\documentclass[aspectratio=169,11pt,xcolor=table]{beamer}
\usepackage{slidestyle}
\ifdefined\notesmode
  \setbeameroption{show only notes}
\else
  \setbeameroption{hide notes}
\fi
\title{Human Reproduction}

\begin{document}

% ================================================================ opening
\begin{frame}
\relax
\vspace*{4mm}
{\usebeamerfont{title}\usebeamercolor[fg]{frametitle}Human Reproduction}\par
\vspace{1mm}{\color{rule}\rule{\textwidth}{0.8pt}}\par\vspace{2mm}
Year 9 Science \textperiodcentered{} exam sprint \textperiodcentered{} 2 hours\par\vspace{4mm}
\begin{tabular}{@{}l@{\hspace{6mm}}l@{}}
0:04 & Loop 1 \textperiodcentered{} Sex cells and the two systems\\
0:26 & Loop 2 \textperiodcentered{} Hormones and the menstrual cycle\\
0:49 & Loop 3 \textperiodcentered{} Fertilisation, chromosomes, twins\\
1:12 & Loop 4 \textperiodcentered{} Pregnancy, birth, fertility\\
1:44 & Mini mock\\
\end{tabular}
\note{%
Teaching line, not a question slide.\par
Plan: opening 0:00, Loop 1 0:04, Loop 2 0:26, Loop 3 0:49, Loop 4 1:12, lose marks and summary 1:42, mini mock 1:44, homework 1:59.\par
Each loop ends with about 6 minutes of writing on the In-class work sheet (Parts A to D). Part E is the mini mock.\par
Say it:\par
reproduction = ree-pruh-DUK-shun\par
hormones = HOR-mohnz\par
menstrual = MEN-stroo-ul\par
fertilisation = fer-tih-ly-ZAY-shun\par
chromosomes = KROH-muh-sohmz\par
pregnancy = PREG-nun-see\par
fertility = fer-TIL-ih-tee}
\end{frame}

\begin{frame}{Do Now}
\relax
\qline{1} Name the male sex cell and the female sex cell.
\qline{2} Which organ makes the female sex cells?
\qline{3} What is the joining of the two sex cells called?
\qline{4} How many chromosomes are in a human body cell?
\qline{5} Which organ passes oxygen from the mother to the foetus?
\fref{Booklet p.2--3}
\note{%
1. sperm (male); egg or ovum (female)\par
2. the ovaries\par
3. fertilisation\par
4. 46 (23 pairs)\par
5. the placenta\par
Diagnostic:\par
1 to 3 wrong: go slowly through Loop 1. 4 wrong: spend longer on the chromosome slide in Loop 3. 5 wrong: spend longer on the placenta slides in Loop 4.\par
Say it:\par
ovum = OH-vum\par
ovaries = OH-vuh-reez\par
placenta = pluh-SEN-tuh\par
foetus = FEE-tus}
\end{frame}

\begin{frame}{The one idea: from two cells to a baby}
\relax
\vspace{3mm}
\FigJourney\par\vspace{5mm}
Every question is about one of these steps.\par\vspace{1mm}
For each step, know the cell, the place and the process.\par\vspace{3mm}
Which step happens in the oviduct?
\fref{Booklet p.3}
\note{%
Teaching line, not a question slide.\par
Q fertilisation happens in the oviduct; the zygote starts to divide there on its way to the uterus\par
Say it:\par
zygote = ZY-goht\par
embryo = EM-bree-oh\par
implantation = im-plan-TAY-shun\par
oviduct = OH-vih-dukt\par
uterus = YOO-tuh-rus}
\end{frame}

% ================================================================ LOOP 1
\begin{frame}{Quick review: Sex cells}
\relax
\twocol{0.52}{\vspace{1mm}\FigSperm[8.6mm]\par\vspace{1mm}\hspace*{6mm}\FigEgg[4.3mm]}{0.44}{%
\qline{1} Name parts A to D and parts P to S.
\qline{2} How many chromosomes are in each nucleus?
\qline{3} Which cell is larger? Which cell can move?}
\fref{Booklet p.4--5 \textperiodcentered{} not to scale}
\note{%
1. A tip of the head (it holds enzymes), B nucleus, C middle piece (it holds mitochondria), D tail; P nucleus, Q cytoplasm, R cell membrane, S jelly layer\par
2. 23 in each\par
3. the egg is much larger; only the sperm can move: it swims with its tail\par
Say it:\par
nucleus = NYOO-klee-us\par
cytoplasm = SY-toh-plaz-um\par
membrane = MEM-brayn\par
enzymes = EN-zymz\par
mitochondria = my-toh-KON-dree-uh}
\end{frame}

\begin{frame}{Adaptations: the feature and how it helps}
\relax
\renewcommand\arraystretch{1.3}
\begin{tabular}{@{}l l l@{}}
\toprule
\textbf{Cell} & \textbf{Feature} & \textbf{How it helps}\\
\midrule
sperm & long tail & \gap[52mm]\\
sperm & many mitochondria & \gap[52mm]\\
sperm & enzymes in the head & \gap[52mm]\\
egg & large cytoplasm & \gap[52mm]\\
egg & membrane changes & \gap[52mm]\\
\bottomrule
\end{tabular}\par\vspace{3mm}
Why does a man make millions of sperm, but a woman releases one egg a month?
\fref{Booklet p.5}
\note{%
Teaching line, not a question slide.\par
Gaps: it swims to the egg; they release energy for swimming; they break through the jelly layer and membrane of the egg; it is a food store for the first cell divisions; it stops any other sperm getting in.\par
Q very few sperm survive the journey to the oviduct, so millions give a better chance that one meets the egg; one egg is enough because it stays inside the body\par
Say it:\par
adaptations = ad-ap-TAY-shunz}
\end{frame}

\begin{frame}{Quick review: Female system, front view}
\relax
\twocol{0.56}{\vspace{1mm}\FigFemaleFront[6.3mm]}{0.40}{%
\qline{1} Name parts A to F.
\qline{2} Where are eggs made?
\qline{3} Where does a foetus develop?
\qline{4} Which part is a ring of muscle?}
\fref{Booklet p.8--9}
\note{%
1. A ovary, B oviduct (fallopian tube), C uterus, D lining of the uterus (endometrium), E cervix, F vagina\par
2. in the ovaries\par
3. in the uterus\par
4. the cervix\par
Say it:\par
ovary = OH-vuh-ree\par
fallopian = fuh-LOH-pee-un\par
endometrium = en-doh-MEE-tree-um\par
cervix = SUR-viks\par
vagina = vuh-JY-nuh}
\end{frame}

\begin{frame}{Female system: the job of each part}
\relax
\renewcommand\arraystretch{1.3}
\begin{tabular}{@{}l l@{}}
\toprule
\textbf{Part} & \textbf{Job}\\
\midrule
ovary & makes \gap[20mm] and two hormones\\
oviduct & carries the egg; \gap[30mm] happens here\\
uterus & the \gap[22mm] develops here\\
lining & gets thick so that an \gap[22mm] can implant\\
cervix & a ring of \gap[22mm] at the base of the uterus\\
vagina & receives \gap[20mm]; the baby is born through it\\
\bottomrule
\end{tabular}
\fref{Booklet p.9}
\note{%
Teaching line, not a question slide.\par
Gaps: eggs; fertilisation; foetus; embryo; muscle; sperm (semen).\par
The two hormones from the ovary are oestrogen and progesterone.\par
Say it:\par
oestrogen = EE-struh-jun\par
progesterone = proh-JES-tuh-rohn\par
semen = SEE-mun}
\end{frame}

\begin{frame}{Quick review: Female system, side view}
\relax
\twocol{0.50}{\vspace{1mm}\FigFemaleSide[72mm]}{0.46}{%
anus \textperiodcentered{} bladder \textperiodcentered{} cervix \textperiodcentered{} fallopian tube \textperiodcentered{} ovary \textperiodcentered{} rectum \textperiodcentered{} urethra \textperiodcentered{} uterus \textperiodcentered{} vagina \textperiodcentered{} vulva
\qline{1} Match the words to parts 1 to 10.
\qline{2} Which four parts are not reproductive organs?}
\fref{Booklet p.8}
\note{%
1. 1 ovary, 2 fallopian tube, 3 bladder, 4 vagina, 5 uterus, 6 rectum, 7 vulva, 8 urethra, 9 anus, 10 cervix\par
2. bladder and urethra (they store and carry urine); rectum and anus (digestive system)\par
Say it:\par
anus = AY-nus\par
rectum = REK-tum\par
urethra = yoo-REE-thruh\par
vulva = VUL-vuh\par
urine = YOOR-in}
\end{frame}

\begin{frame}{Quick review: Male system}
\relax
\twocol{0.50}{\vspace{1mm}\FigMaleFront[60mm]}{0.46}{%
\qline{1} Name parts A to H.
\qline{2} Where are sperm made?
\qline{3} Which tube carries both urine and semen?}
\fref{Booklet p.10}
\note{%
1. A bladder, B seminal vesicle, C prostate gland, D sperm duct (vas deferens), E urethra, F epididymis, G testis, H scrotum\par
2. in the testes\par
3. the urethra\par
Say it:\par
seminal = SEM-ih-nul\par
vesicle = VES-ih-kul\par
prostate = PROS-tayt\par
vas deferens = vas-DEF-uh-renz\par
epididymis = ep-ih-DID-ih-mis\par
testis = TES-tis\par
testes = TES-teez\par
scrotum = SKROH-tum}
\end{frame}

\begin{frame}{The path of a sperm cell}
\relax
Put the cards in order.\par\vspace{1mm}
\qline{A} The urethra carries semen out through the penis.
\qline{B} Sperm are made in a testis.
\qline{C} Sperm travel along the sperm duct.
\qline{D} Sperm are stored in the epididymis.
\qline{E} Glands add fluid: sperm + fluid = \gap[24mm].
\par\vspace{3mm}
Why are the testes held outside the body?
\fref{Booklet p.11}
\note{%
Teaching line, not a question slide.\par
Order: B, D, C, E, A. Gap: semen.\par
Q sperm are made best at a temperature a little below body temperature; in the scrotum the testes stay cooler\par
Say it:\par
penis = PEE-nis}
\end{frame}

\begin{frame}{Male system: the job of each part}
\relax
\renewcommand\arraystretch{1.3}
\begin{tabular}{@{}l l@{}}
\toprule
\textbf{Part} & \textbf{Job}\\
\midrule
scrotum & holds the testes outside the body, so they stay \gap[20mm]\\
testes & make \gap[18mm] and the hormone \gap[32mm]\\
epididymis & \gap[20mm] the sperm\\
sperm duct & carries sperm to the \gap[24mm]\\
glands & add \gap[18mm] that feeds the sperm\\
penis & puts semen into the \gap[22mm]\\
\bottomrule
\end{tabular}
\fref{Booklet p.11}
\note{%
Teaching line, not a question slide.\par
Gaps: cooler; sperm, testosterone; stores; urethra; fluid; vagina.\par
The glands are the seminal vesicles and the prostate gland.\par
Say it:\par
testosterone = tes-TOS-tuh-rohn}
\end{frame}

\begin{frame}{Spot the error (1)}
\relax
Describe where sex cells are made and where a foetus develops.\par\vspace{2mm}
\samcard{146mm}{\flawed{Eggs are made in the uterus and travel down the oviduct. The cervix is the place where the foetus develops. In a man, sperm are made in the testes and stored in the scrotum.}}\par\vspace{2mm}
\textbf{Sam's answer has 3 errors.} Find each one, fix it, and say why it loses marks.
\fref{Booklet p.8--11}
\note{%
1. eggs are made in the ovaries, not in the uterus (\Lref{parts})\par
2. the foetus develops in the uterus; the cervix is the ring of muscle at its base (\Lref{parts})\par
3. sperm are stored in the epididymis; the scrotum is the bag of skin that holds the testes (\Lref{parts})\par
Short solution:\par
Eggs are made in the ovaries and travel down the oviduct. The uterus is the place where the foetus develops. In a man, sperm are made in the testes and stored in the epididymis.}
\end{frame}

\begin{frame}{Practice}
\relax
\qline{Q1} State two differences between a sperm cell and an egg cell.\slv{A}
\qline{Q2} Explain why the testes are held outside the body in the scrotum.\slv{M}
\qline{Q3} A woman's oviducts are both blocked. She still releases an egg each month. Explain why she cannot become pregnant naturally.\slv{E}
\fref{Booklet p.4--5, p.9--11}
\note{%
Q1 any two: a sperm is much smaller; a sperm has a tail and can swim, an egg cannot move; an egg has a large food store in its cytoplasm; millions of sperm are made, one egg is released each month\par
Q2 sperm are made best at a temperature a little below body temperature; outside the body the testes stay cooler\par
Q3 fertilisation happens in the oviduct; the sperm cannot reach the egg and the egg cannot reach the uterus, so no zygote forms and nothing can implant}
\end{frame}

\begin{frame}{Classwork}
\relax
\vspace{4mm}
\emphline{In-class work \textperiodcentered{} Part A}\par\vspace{4mm}
\qline{1} Label the female reproductive system.
\qline{2} Complete the sentences about the male system.
\qline{3} Explain one adaptation of a sperm cell.
\par\vspace{5mm}
Write in full sentences. Use the science words.
\note{%
Teaching line, not a question slide.\par
About 6 minutes, 0:20 to 0:26. The answers are in inclass\_answers.pdf.}
\end{frame}

% ================================================================ LOOP 2
\begin{frame}{Puberty: hormones switch it on}
\relax
Hormones are chemical \gap[30mm], carried in the \gap[18mm].\par\vspace{1mm}
\renewcommand\arraystretch{1.3}
\begin{tabular}{@{}l l l@{}}
\toprule
 & \textbf{Boys} & \textbf{Girls}\\
\midrule
hormone & \gap[32mm] & \gap[28mm]\\
made in & testes & ovaries\\
changes & voice deepens & \gap[22mm] start\\
also & sperm are made & breasts grow\\
\bottomrule
\end{tabular}\par\vspace{3mm}
True or false? Girls usually start puberty before boys.
\fref{Booklet p.6--7, p.13}
\note{%
Teaching line, not a question slide.\par
Gaps: messengers; blood; testosterone; oestrogen; periods.\par
Q true\par
Both sexes: a growth spurt, and pubic and underarm hair.\par
Say it:\par
puberty = PYOO-buh-tee}
\end{frame}

\begin{frame}{The menstrual cycle: stages 1 and 2}
\relax
\FigCycleBar[4.6mm]\par\vspace{1mm}
\qline{1} Days 1 to 5: the lining \gap[30mm] and leaves through the vagina. This is a \gap[20mm].
\qline{2} Days 5 to 13: the lining \gap[26mm] again. An egg matures in an \gap[20mm].
\fref{Booklet p.14--15}
\note{%
Teaching line, not a question slide.\par
Gaps: breaks down; period (menstruation); builds up; ovary.\par
Stage names: 1 menstrual, 2 follicular. In stage 2 a follicle in the ovary grows; FSH (follicle stimulating hormone) makes the egg inside it mature.\par
Say it:\par
menstruation = men-stroo-AY-shun\par
follicular = fuh-LIK-yoo-luh\par
follicle = FOL-ih-kul\par
stimulating = STIM-yoo-lay-ting}
\end{frame}

\begin{frame}{The menstrual cycle: stages 3 and 4}
\relax
\FigCycleBar[4.6mm]\par\vspace{1mm}
\qline{3} About day 14: an egg is \gap[24mm] from an ovary. This is \gap[26mm].
\qline{4} Days 15 to 28: the lining stays \gap[18mm], ready for a fertilised egg.
\par\vspace{1mm}
When is a woman most likely to become pregnant?
\fref{Booklet p.14--16}
\note{%
Teaching line, not a question slide.\par
Gaps: released; ovulation; thick.\par
Q around day 14: on the day of ovulation and the two days before and after it\par
Stage names: 3 ovulation, 4 luteal. With no fertilised egg the lining breaks down and the cycle starts again.\par
Say it:\par
ovulation = ov-yoo-LAY-shun\par
luteal = LOO-tee-ul}
\end{frame}

\begin{frame}{Four hormones, four jobs}
\relax
\renewcommand\arraystretch{1.3}
\begin{tabular}{@{}l l l@{}}
\toprule
\textbf{Hormone} & \textbf{Made in} & \textbf{Job}\\
\midrule
FSH & pituitary gland & makes an egg \gap[20mm] in the ovary\\
LH & \gap[34mm] & triggers \gap[24mm] on day 14\\
oestrogen & ovary & \gap[26mm] the lining of the uterus\\
progesterone & \gap[34mm] & keeps the lining \gap[18mm]\\
\bottomrule
\end{tabular}\par\vspace{3mm}
FSH = follicle stimulating hormone. LH = luteinising hormone.\par\vspace{1mm}
Which two hormones peak at ovulation?
\fref{Booklet p.13--14}
\note{%
Teaching line, not a question slide.\par
Gaps: mature; pituitary gland, ovulation; builds up; ovary, thick.\par
Q FSH and LH\par
Say it:\par
pituitary = pih-TYOO-ih-tree\par
luteinising = LOO-tee-in-eye-zing}
\end{frame}

\begin{frame}{Read the graph: ovarian hormones}
\relax
\twocol{0.54}{\vspace{1mm}\FigHormones[2.55mm]{ov}}{0.43}{%
\qline{1} Which hormone peaks just before day 14?
\qline{2} Which hormone is highest at about day 21?
\qline{3} Both levels fall by day 28. What happens to the lining?}
\fref{Booklet p.14}
\note{%
Teaching line, not a question slide.\par
1. oestrogen\par
2. progesterone\par
3. the lining breaks down and is lost: a period starts}
\end{frame}

\begin{frame}{Read the graph: pituitary hormones}
\relax
\twocol{0.54}{\vspace{1mm}\FigHormones[2.55mm]{pit}}{0.43}{%
\qline{1} On which day do LH and FSH peak?
\qline{2} What does the LH peak cause?
\qline{3} Describe the level of LH from day 1 to day 12.}
\fref{Booklet p.14}
\note{%
Teaching line, not a question slide.\par
1. day 14\par
2. ovulation: an egg is released from an ovary\par
3. it stays low and almost constant}
\end{frame}

\begin{frame}{No fertilised egg: write it in this order}
\relax
\qline{1} The egg is not fertilised, so no embryo \gap[24mm] in the lining.
\qline{2} The level of \gap[34mm] falls.
\qline{3} The lining of the uterus \gap[32mm].
\qline{4} Blood and lining leave through the \gap[22mm]: a period.
\par\vspace{2mm}
{\bfseries A}: lines 3--4 \quad {\bfseries M}: links the hormone to the lining \quad {\bfseries E}: the whole chain\par\vspace{2mm}
What is different if the egg is fertilised?
\fref{Booklet p.9, p.15}
\note{%
Teaching line, not a question slide.\par
Gaps: implants; progesterone; breaks down; vagina.\par
Q the embryo implants, progesterone stays high and the lining stays thick, so there is no period}
\end{frame}

\begin{frame}{Spot the error (2)}
\relax
Describe ovulation and name one hormone that changes the lining of the uterus.\par\vspace{2mm}
\samcard{146mm}{\flawed{Ovulation is when the lining of the uterus is lost. An egg is released on about day 28 of the cycle. Oestrogen is made in the pituitary gland and makes the lining thicker.}}\par\vspace{2mm}
\textbf{Sam's answer has 3 errors.} Find each one, fix it, and say why it loses marks.
\fref{Booklet p.13--16}
\note{%
1. ovulation is the release of an egg from an ovary; losing the lining is menstruation (\Lref{ovul})\par
2. the egg is released on about day 14, not day 28 (\Lref{ovul})\par
3. oestrogen is made in the ovaries, not in the pituitary gland (\Lref{horm})\par
Short solution:\par
Ovulation is the release of an egg from an ovary, on about day 14 of the cycle. Oestrogen is made in the ovaries and makes the lining thicker.}
\end{frame}

\begin{frame}{Practice}
\relax
\qline{Q1} A woman has a 28-day cycle. On which days is she most likely to become pregnant?\slv{A}
\qline{Q2} Describe how the lining of the uterus changes between day 5 and day 14. Name the hormone that causes this.\slv{M}
\qline{Q3} The level of progesterone stays high in a pregnant woman. Explain why this is important.\slv{E}
\fref{Booklet p.13--16}
\note{%
Q1 about days 12 to 16: ovulation is on about day 14, and she is fertile for two days before and after it\par
Q2 the lining builds up and becomes thicker, with more blood vessels; oestrogen\par
Q3 progesterone keeps the lining of the uterus thick, so the lining is not lost; the embryo stays implanted and the placenta can form, so there are no periods during pregnancy}
\end{frame}

\begin{frame}{Classwork}
\relax
\vspace{4mm}
\emphline{In-class work \textperiodcentered{} Part B}\par\vspace{4mm}
\qline{1} Name the four stages of the cycle.
\qline{2} Match each hormone to its job.
\qline{3} Read a table of hormone levels and explain it.
\par\vspace{5mm}
For data: describe the trend first, with numbers.
\note{%
Teaching line, not a question slide.\par
About 6 minutes, 0:43 to 0:49. The answers are in inclass\_answers.pdf.}
\end{frame}

% ================================================================ LOOP 3
\begin{frame}{From intercourse to fertilisation}
\relax
Put the cards in order.\par\vspace{1mm}
\qline{A} Sperm swim through the cervix and uterus into the oviducts.
\qline{B} The penis fills with blood and becomes erect.
\qline{C} One sperm enters the egg, and the two nuclei fuse.
\qline{D} Semen is released into the vagina. This is ejaculation.
\qline{E} The penis is put into the vagina.
\par\vspace{3mm}
Which card is fertilisation? Where does it happen?
\fref{Booklet p.17}
\note{%
Teaching line, not a question slide.\par
Order: B, E, D, A, C.\par
Q card C is fertilisation; it happens in the oviduct\par
Say it:\par
intercourse = IN-ter-kors\par
ejaculation = ih-jak-yoo-LAY-shun}
\end{frame}

\begin{frame}{Zygote, embryo, foetus}
\relax
\qline{1} The sperm nucleus fuses with the egg nucleus. The new cell is a \gap[22mm].
\qline{2} It divides many times as it moves along the oviduct. The ball of cells is an \gap[22mm].
\qline{3} The ball of cells sinks into the lining of the uterus. This is \gap[32mm].
\qline{4} After about 8 weeks it has all its organs. Now it is a \gap[22mm].
\par\vspace{3mm}
What stops a second sperm from entering the egg?
\fref{Booklet p.17, p.21}
\note{%
Teaching line, not a question slide.\par
Gaps: zygote; embryo; implantation; foetus.\par
Q the membrane of the egg changes as soon as one sperm has entered}
\end{frame}

\begin{frame}{Chromosomes: count them}
\relax
\vspace{2mm}
\FigChromo{\gap}{\gap}{\gap}{\gap}\par\vspace{4mm}
Chromosomes carry the genes. They are in the \gap[24mm] of every cell.\par\vspace{2mm}
Why must a sex cell have only half the number?
\fref{Booklet p.18}
\note{%
Teaching line, not a question slide.\par
Gaps: sperm 23; egg 23; zygote 46; every body cell 46 (23 pairs); nucleus.\par
Q so that the zygote has the full 46 and not 92 when the two sex cells join; half comes from each parent\par
Say it:\par
genes = JEENZ}
\end{frame}

\begin{frame}{Boy or girl?}
\relax
\twocol{0.43}{\vspace{2mm}\FigSexSquare{\gap}{\gap}{\gap}{\gap}}{0.54}{%
Female cells have XX. Male cells have XY.\par\vspace{1mm}
Every egg carries \gap. A sperm carries \gap\ or \gap.
\qline{1} Complete the square.
\qline{2} What is the chance of a boy?
\qline{3} Which parent's sex cell decides?}
\fref{Booklet p.20}
\note{%
Teaching line, not a question slide.\par
Gaps: X; X or Y.\par
1. top row XX, XX; bottom row XY, XY\par
2. 50\,\% (2 of the 4 boxes)\par
3. the father's: a sperm with X gives a girl, a sperm with Y gives a boy}
\end{frame}

\begin{frame}{Twins}
\relax
\twocol{0.40}{\vspace{2mm}\FigTwins[5.2mm]}{0.57}{%
Each row ends with two embryos.
\qline{1} Which row shows identical twins?
\qline{2} Row X: \gap[12mm] egg, \gap[12mm] sperm. The zygote \gap[18mm].
\qline{3} Row Y: \gap[12mm] eggs, \gap[12mm] sperm.
\qline{4} Which twins can be a boy and a girl?}
\fref{Booklet p.19}
\note{%
Teaching line, not a question slide.\par
1. row X\par
2. one, one; splits into two\par
3. two, two\par
4. row Y, the non-identical (fraternal) twins: they have different genes; identical twins have the same genes, so they are always the same sex\par
Say it:\par
identical = eye-DEN-tih-kul\par
fraternal = fruh-TUR-nul}
\end{frame}

\begin{frame}{Workbook: true or false?}
\relax
\qline{a} Identical twins happen when a fertilised egg cell divides in two, but the two cells split apart.
\qline{b} Identical twins will always be different sexes.
\qline{c} Non-identical twins happen when two egg cells are fertilised.
\qline{d} Non-identical twins can be the same sex or different sexes.
\qline{e} Identical twins happen when two egg cells are fertilised.
\qline{f} Non-identical twins happen when egg cells turn into embryos on their own.
\fref{Booklet p.19}
\note{%
(a) true\par
(b) false: they are always the same sex\par
(c) true\par
(d) true\par
(e) false: that gives non-identical twins\par
(f) false: each egg must be fertilised by a sperm}
\end{frame}

\begin{frame}{Workbook: fertilisation}
\relax
\qline{Q1} Where does fertilisation occur?
\qline{Q2} Only one sperm cell joins with the egg cell. What do you think happens to all the other sperms?
\qline{Q3} Humans have internal fertilisation as a result of moving onto the land in our evolutionary history. Why do you think this is an advantage on the land?
\fref{Booklet p.18}
\note{%
Q1 in the oviduct (fallopian tube)\par
Q2 they die and are broken down\par
Q3 on land the sex cells would dry out; inside the body the sperm can swim in fluid, the sex cells are protected, and a sperm is much more likely to meet the egg\par
Say it:\par
evolutionary = ee-vuh-LOO-shun-ree}
\end{frame}

\begin{frame}{A boy next time? Write it in this order}
\relax
A couple have three girls. What is the chance that their fourth child is a boy?\par\vspace{1mm}
\qline{1} Every egg carries an \gap\ chromosome.
\qline{2} Half of the sperm carry X and half carry \gap.
\qline{3} So at every fertilisation the chance of a boy is \gap[18mm].
\qline{4} The children born before do not change the \gap[20mm].
\par\vspace{2mm}
{\bfseries A}: the number \quad {\bfseries M}: lines 1--2 \quad {\bfseries E}: line 4 as well
\fref{Booklet p.20}
\note{%
Teaching line, not a question slide.\par
Gaps: X; Y; 50\,\% (1 in 2); chance.\par
Each fertilisation is a new event with a new sperm cell.}
\end{frame}

\begin{frame}{Spot the error (3)}
\relax
Explain where a baby's chromosomes come from, and how twins can be identical.\par\vspace{2mm}
\samcard{146mm}{\flawed{A sperm has 46 chromosomes and an egg has 46, so the baby gets all of them. The mother's egg decides if the baby is a boy or a girl. Identical twins form when two eggs are fertilised by two sperm.}}\par\vspace{2mm}
\textbf{Sam's answer has 3 errors.} Find each one, fix it, and say why it loses marks.
\fref{Booklet p.18--20}
\note{%
1. a sperm and an egg have 23 chromosomes each, so the zygote has 46 (\Lref{chromo})\par
2. the sperm decides the sex: it carries X or Y, and every egg carries X (\Lref{sex})\par
3. identical twins come from one fertilised egg that splits; two eggs give non-identical twins (\Lref{twins})\par
Short solution:\par
A sperm has 23 chromosomes and an egg has 23, so the zygote has 46. The father's sperm decides if the baby is a boy or a girl. Identical twins form when one fertilised egg splits into two.}
\end{frame}

\begin{frame}{Practice}
\relax
\qline{Q1} State where fertilisation happens and name the cell that is formed.\slv{A}
\qline{Q2} A boy and a girl are twins. Are they identical twins? Explain.\slv{M}
\qline{Q3} Explain why a child has features of both parents but is not identical to either of them.\slv{E}
\fref{Booklet p.3, p.18--20}
\note{%
Q1 in the oviduct; a zygote\par
Q2 no: identical twins come from one zygote, so they have the same chromosomes and are the same sex; a boy and a girl must come from two eggs and two sperm\par
Q3 the child gets 23 chromosomes from the mother's egg and 23 from the father's sperm, so its genes are a mixture of both; the mixture is a new combination, different from each parent}
\end{frame}

\begin{frame}{Classwork}
\relax
\vspace{4mm}
\emphline{In-class work \textperiodcentered{} Part C}\par\vspace{4mm}
\qline{1} Put the stages from sex cell to foetus in order.
\qline{2} Work out the sex of a baby from its chromosomes.
\qline{3} Explain how two kinds of twins form.
\par\vspace{5mm}
Give the number of chromosomes each time you name a cell.
\note{%
Teaching line, not a question slide.\par
About 6 minutes, 1:06 to 1:12. The answers are in inclass\_answers.pdf.}
\end{frame}

% ================================================================ LOOP 4
\begin{frame}{Quick review: The foetus in the uterus}
\relax
\twocol{0.50}{\vspace{1mm}\FigWomb[4.7mm]}{0.46}{%
\qline{1} Name parts A to F.
\qline{2} Which part cushions the foetus?
\qline{3} Which part joins the foetus to the placenta?}
\fref{Booklet p.21, p.25}
\note{%
1. A placenta, B umbilical cord, C amniotic fluid, D amniotic sac, E uterus wall, F cervix\par
2. the amniotic fluid\par
3. the umbilical cord\par
Say it:\par
umbilical = um-BIL-ih-kul\par
amniotic = am-nee-OT-ik}
\end{frame}

\begin{frame}{The placenta: what crosses?}
\relax
\vspace{1mm}
\hbox to\textwidth{\hfil\FigExchange{\gap[18mm]}{\gap[18mm]}{\gap[18mm]}{\gap[18mm]}{\gap[18mm]}\hfil}\par\vspace{3mm}
Down: three things the foetus needs. \quad Up: two waste products.\par\vspace{2mm}
The two bloods never \gap[16mm]. Substances cross by \gap[28mm].
\fref{Booklet p.22--24}
\note{%
Teaching line, not a question slide.\par
Down: oxygen; glucose (food, nutrients); antibodies. Up: carbon dioxide; urea.\par
Gaps: mix; diffusion.\par
Say it:\par
villi = VIL-eye\par
glucose = GLOO-kohz\par
antibodies = AN-tee-bod-eez\par
urea = yoo-REE-uh\par
diffusion = dih-FYOO-zhun}
\end{frame}

\begin{frame}{Workbook: the role of the placenta}
\relax
The placenta develops from foetal tissues and allows the \gap[26mm] of materials between mother and foetus. It has a large number of blood vessels that do not touch. This ensures there is no \gap[22mm] of the maternal and foetal blood.
\qline{a} Give three reasons why it is important that the maternal and foetal blood do not mix.
\qline{b} Explain why the embryo needs oxygen and glucose.
\qline{c} Which substances diffuse from the foetus' blood into the mother's blood and why is it important that they are removed?
\fref{Booklet p.26}
\note{%
(a) the mother and the foetus can have different blood types; the mother's blood pressure is higher and would damage the foetus's small blood vessels; keeping the bloods apart stops many bacteria and harmful substances reaching the foetus\par
(b) for respiration, which releases the energy its cells need to divide and grow\par
(c) carbon dioxide and urea; they are waste products and would poison the foetus if they built up\par
Short solution:\par
Blanks: exchange (diffusion); mixing.\par
Say it:\par
foetal = FEE-tul\par
maternal = muh-TUR-nul\par
respiration = res-pih-RAY-shun}
\end{frame}

\begin{frame}{Oxygen to the foetus: write it in this order}
\relax
\qline{1} The mother's blood has a \gap[20mm] concentration of oxygen than the foetus's blood.
\qline{2} Oxygen \gap[24mm] from high to low concentration, across the thin wall.
\qline{3} The villi give a large \gap[32mm], so diffusion is fast.
\qline{4} The foetus's blood carries the oxygen away in the umbilical \gap[16mm].
\par\vspace{1mm}
{\bfseries A}: line 2 \quad {\bfseries M}: lines 1--2 \quad {\bfseries E}: lines 3--4 as well\par\vspace{1mm}
Now tell the same story for carbon dioxide.
\fref{Booklet p.22--24}
\note{%
Teaching line, not a question slide.\par
Gaps: higher; diffuses; surface area; vein.\par
Q the foetus's blood, arriving in the umbilical artery, has a higher concentration of carbon dioxide than the mother's blood, so carbon dioxide diffuses across the villi into the mother's blood, which carries it away\par
Say it:\par
concentration = kon-sun-TRAY-shun\par
artery = AR-tuh-ree}
\end{frame}

\begin{frame}{Harm: what else crosses the placenta?}
\relax
\renewcommand\arraystretch{1.3}
\begin{tabular}{@{}l l@{}}
\toprule
\textbf{From the mother} & \textbf{Effect on the foetus}\\
\midrule
cigarette smoke & less \gap[20mm] reaches it; the baby is smaller\\
alcohol & damages the developing \gap[18mm]\\
drugs & the baby can be born addicted\\
rubella virus & deafness; heart and eye defects\\
\bottomrule
\end{tabular}\par\vspace{3mm}
Why can even a small amount do a lot of damage?
\fref{Booklet p.29}
\note{%
Teaching line, not a question slide.\par
Gaps: oxygen; brain.\par
Q the foetus is very small and is growing very quickly\par
Cigarette smoke: carbon monoxide takes the place of oxygen in the mother's blood; nicotine also crosses the placenta.\par
Say it:\par
rubella = roo-BEL-uh\par
addicted = uh-DIK-tid\par
monoxide = mon-OK-syd\par
nicotine = NIK-uh-teen}
\end{frame}

\begin{frame}{Pregnancy in numbers}
\relax
\qline{1} Pregnancy (gestation) lasts about \gap\ weeks, counted from the first day of the last period.
\qline{2} The embryo is called a foetus from about week \gap.
\qline{3} Losing a foetus before \gap\ weeks is a miscarriage.
\qline{4} A baby born before \gap\ weeks is premature.
\par\vspace{3mm}
A baby is born at 35 weeks. Which word describes this birth?
\fref{Booklet p.27--28, p.37}
\note{%
Teaching line, not a question slide.\par
Gaps: 40; 8 (the glossary says after 8 weeks, page 28 says nine); 20; 38.\par
Q premature\par
After 20 weeks the loss of a baby is called a stillbirth.\par
Say it:\par
gestation = jes-TAY-shun\par
miscarriage = MIS-ka-rij\par
premature = PREM-uh-chur}
\end{frame}

\begin{frame}{Birth: three stages}
\relax
Put the cards in order. Then name each stage.\par\vspace{1mm}
\qline{A} The placenta is pushed out.
\qline{B} Muscles of the uterus contract. The cervix opens (dilates). The waters break.
\qline{C} The baby is pushed out through the vagina. The cord is tied and cut.
\par\vspace{2mm}
Stage 1 \gap[24mm] \quad Stage 2 \gap[24mm] \quad Stage 3 \gap[24mm]\par\vspace{3mm}
What are ``the waters''?
\fref{Booklet p.30--31}
\note{%
Teaching line, not a question slide.\par
Order: B, C, A. Stage 1 labour; stage 2 delivery; stage 3 afterbirth.\par
Q the amniotic fluid\par
Say it:\par
dilates = dy-LAYTS\par
labour = LAY-bur}
\end{frame}

\begin{frame}{Workbook: fertility problems and solutions}
\relax
\renewcommand\arraystretch{1.15}
\begin{tabular}{@{}p{62mm} p{92mm}@{}}
\toprule
\textbf{Problems} & \textbf{Solutions}\\
\midrule
Low sperm count & Injecting sperm cell into egg cell, then IVF\\
Oviducts too narrow & \gap[40mm]\\
Sperm ducts blocked & \gap[40mm]\\
Lining stops implantation & \gap[40mm]\\
Few ripe eggs released & \gap[40mm]\\
Sperm do not swim well & \gap[40mm]\\
\bottomrule
\end{tabular}
\fref{Booklet p.31--32}
\note{%
1. doctors try to widen the oviducts; or IVF\par
2. doctors try to widen the sperm ducts\par
3. IVF\par
4. IVF\par
5. inject a sperm cell into an egg cell with a fine needle\par
Short solution:\par
These are the answers in the paragraphs on page 31. The numbers 1 to 5 are the five empty rows, from the top.}
\end{frame}

\begin{frame}{IVF: fertilisation in a dish}
\relax
Put the cards in order.\par\vspace{1mm}
\qline{A} Eggs are collected from the woman's ovaries.
\qline{B} An embryo is put into the uterus, where it may implant.
\qline{C} The woman is given hormones, so several eggs mature.
\qline{D} The eggs are mixed with sperm in a dish. Some are fertilised.
\qline{E} The embryos grow for about two days.
\par\vspace{3mm}
Why are twins more common after IVF?
\fref{Booklet p.33--34}
\note{%
Teaching line, not a question slide.\par
Order: C, A, D, E, B.\par
Q often more than one embryo is put into the uterus, and each one can implant\par
In vitro means in glass: fertilisation happens outside the body.\par
Say it:\par
vitro = VEE-troh}
\end{frame}

\begin{frame}{Workbook: question 3, describe}
\relax
\twocol{0.56}{\vspace{0mm}\FigMultiBirths[1.22mm]}{0.42}{%
\qline[9mm]{(a)} Describe what has happened to the number of multiple births in women aged 40--44.
\qline[9mm]{(c)} Describe what happened to the number of multiple births in women aged 20--24.}
\fref{Booklet p.35--36 \textperiodcentered{} graph re-drawn}
\note{%
Q3 (a) from 1960 to about 1990 it stayed between about 11\,\% and 15\,\%; after 1990 it rose steadily to about 25\,\% in 2010\par
Q3 (c) it stayed almost the same, at about 9\,\% to 10\,\%, from 1960 to 2010\par
Short solution:\par
A multiple birth is two or more babies born from one pregnancy.}
\end{frame}

\begin{frame}{Workbook: suggest and predict}
\relax
Headlines: ``Fertility treatment shown to increase chances of multiple birth'' \textperiodcentered{} ``Older women more likely to have problems becoming pregnant'' \textperiodcentered{} ``Fertility treatment becoming more common''
\qline{Q3} (b) Suggest a hypothesis for why this has been happening.
\qline{Q3} (d) Suggest a hypothesis for why this pattern is not like that for women aged 40--44.
\qline{Q4} (a) Predict the percentage of multiple births last year for women aged 40--44.
\fref{Booklet p.35--36}
\note{%
Q3 (b) fertility treatment has become more common; older women are more likely to need it; fertility treatment increases the chance of a multiple birth\par
Q3 (d) women aged 20 to 24 rarely have problems becoming pregnant, so few of them have fertility treatment\par
Q4 (a) about 30\,\% or a little more: any value above 25\,\% that continues the rising line\par
Q4 (b) I extended the trend line from 1990 to 2010 forward to last year\par
Say it:\par
hypothesis = hy-POTH-uh-sis}
\end{frame}

\begin{frame}{Spot the error (4)}
\relax
Describe how a foetus gets oxygen and how it is protected in the uterus.\par\vspace{2mm}
\samcard{146mm}{\flawed{The mother's blood flows through the umbilical cord into the foetus. Oxygen moves by diffusion from a low concentration to a high concentration. The amniotic fluid gives the foetus its food and protects it from bumps.}}\par\vspace{2mm}
\textbf{Sam's answer has 3 errors.} Find each one, fix it, and say why it loses marks.
\fref{Booklet p.21--25}
\note{%
1. the two bloods never mix: only the foetus's blood flows in the umbilical cord (\Lref{nomix})\par
2. diffusion is from a high concentration to a low concentration (\Lref{diff})\par
3. the amniotic fluid does not feed the foetus: food comes from the mother's blood through the placenta (\Lref{fluid})\par
Short solution:\par
Oxygen diffuses from the mother's blood into the foetus's blood in the placenta, from a high concentration to a low concentration. The foetus's blood carries it along the umbilical cord. The amniotic fluid protects the foetus from bumps.}
\end{frame}

\begin{frame}{Practice}
\relax
\qline{Q1} Name two substances that pass from the foetus's blood to the mother's blood.\slv{A}
\qline{Q2} Explain why a pregnant woman who smokes is likely to have a smaller baby.\slv{M}
\qline{Q3} IVF clinics now usually put one embryo into the uterus, not three. Suggest why, and explain how this changes the number of twins and triplets.\slv{E}
\fref{Booklet p.22--24, p.29, p.34}
\note{%
Q1 carbon dioxide and urea\par
Q2 carbon monoxide from the smoke is carried in the mother's blood, so less oxygen crosses the placenta; the foetus's cells respire less and release less energy, so it grows more slowly\par
Q3 a multiple birth is riskier for the mother and the babies, who are often premature and small; with one embryo only one can implant, so there are fewer twins and triplets}
\end{frame}

\begin{frame}{Classwork}
\relax
\vspace{4mm}
\emphline{In-class work \textperiodcentered{} Part D}\par\vspace{4mm}
\qline{1} Label the foetus in the uterus.
\qline{2} Explain how a waste product leaves the foetus.
\qline{3} Describe and explain a table about the foetus.
\par\vspace{5mm}
Name the substance, the direction and the process.
\note{%
Teaching line, not a question slide.\par
About 6 minutes, 1:36 to 1:42. The answers are in inclass\_answers.pdf.}
\end{frame}

% ================================================================ closing
\begin{frame}{Ways to lose marks}
\relax
Say the correct version of each one.\par\vspace{1mm}
\qline{1} ``A sperm cell has 46 chromosomes.''
\qline{2} ``Fertilisation happens in the uterus.''
\qline{3} ``Ovulation is when the lining is lost.''
\qline{4} ``The mother's blood flows into the foetus.''
\qline{5} ``A sperm has a tail.'' (asked to explain an adaptation)
\qline{6} ``It went up because of IVF.'' (asked to describe a graph)
\note{%
Teaching line, not a question slide.\par
1. 23: every sex cell has one set; the zygote and body cells have 46 (\Lref{chromo})\par
2. in the oviduct; the embryo implants in the uterus later (\Lref{fert})\par
3. ovulation is the release of an egg; losing the lining is menstruation (\Lref{ovul})\par
4. the bloods never mix; oxygen and food diffuse across the placenta (\Lref{nomix})\par
5. add how it helps: so that it can swim to the egg (\Lref{adapt})\par
6. describe needs the trend and numbers from the graph; the reason is for suggest or explain (\Lref{graph})}
\end{frame}

\begin{frame}{Summary}
\relax
\vspace{2mm}
{\sSum
\qline{1} Sex cells have 23 chromosomes; fertilisation in the oviduct makes 46.
\qline{2} Cycle: period, lining builds up, ovulation on day 14, lining kept thick.
\qline{3} Placenta: diffusion from high to low; the two bloods never mix.
\qline{4} Explain = fact + because. Describe a graph = trend + numbers.\par}
\note{%
Teaching line, not a question slide.\par
Ask her to say each line with the book closed, then give one example question for each.}
\end{frame}

\begin{frame}{Mini mock}
\relax
\vspace{4mm}
\emphline{In-class work \textperiodcentered{} Part E}\par\vspace{4mm}
12 minutes of writing, no notes, no booklet.\par\vspace{2mm}
Then we mark it together: A, M or E for each question.\par\vspace{5mm}
Read the command word first: state, describe, explain, suggest.
\note{%
Teaching line, not a question slide.\par
1:44 to 1:56 writing, 1:56 to 1:59 marking. The answers are in inclass\_answers.pdf.}
\end{frame}

\begin{frame}{Homework}
\relax
\vspace{2mm}
\emphline{Homework: two parts, about 40 minutes each}\par\vspace{3mm}
\qline{1} Part 1: sex cells, the two systems, the menstrual cycle.
\qline{2} Part 2: fertilisation, chromosomes, pregnancy, fertility.
\par\vspace{3mm}
Do the two parts on two different days.
\note{%
Teaching line, not a question slide.\par
The answers are in homework\_answers.pdf.}
\end{frame}

\end{document}
